% Reviewed translation: PDF pages 256--260 / printed pages 242--246.
% The original scan is authoritative; mathematical tokens were checked against it.
\MathBookSourcePage{256}
\subsection{$l$ 次三角形有限元逼近}
\label{subsec:incompressible-hhj-triangular-approximation}

现在逼近问题 \eqref{eq:incompressible-hhj-formulation}.
为简化讨论, 假设 $\mathcal T_h$ 由三角形组成,
但当前方法很容易推广到 $\mathcal T_h$ 还含四边形的情形.
首先选择有限元空间. 像通常一样, 对某个整数 $l\geq1$, 取
\begin{equation}\label{eq:incompressible-hhj-stream-fe-space}
\Theta_h
=\{\theta_h\in C^0(\overline\Omega);
\ \theta_h|_\kappa\in P_l
\ \forall\kappa\in\mathcal T_h\},
\qquad
\Phi_h=\Theta_h\cap\Phi.
\end{equation}
但对张量而言, 由于 $\mathcal X$ 中的张量不必整体连续,
可用自由度更少的空间 $\mathcal X_h$ 来逼近:
\begin{equation}\label{eq:incompressible-hhj-tensor-fe-space}
\mathcal X_h
=\{\boldsymbol\tau_h\in\mathcal X;
\ \tau_{ij}|_\kappa\in P_{l-1}
\ \forall\kappa\in\mathcal T_h,
\ 1\leq i,j\leq2\}.
\end{equation}
用这些空间, 问题 \eqref{eq:incompressible-hhj-formulation}
离散化为问题 $(Q_h)$:

求 $(\boldsymbol\sigma_h,\psi_h)
\in\mathcal X_h\times\Phi_h$, 使
\begin{equation}\label{eq:incompressible-hhj-discrete-problem}
\left\{
\begin{aligned}
\nu\sum_{\kappa\in\mathcal T_h}
\int_\kappa\frac{\partial^2\phi_h}{\partial x_i\partial x_j}
(\sigma_h)_{ij}\,dx
-\nu\int_{\Gamma_h}M_n(\boldsymbol\sigma_h)
S(\partial\phi_h/\partial n)\,ds
&=(\mathbf f,\operatorname{curl}\phi_h)
&&\forall\phi_h\in\Phi_h,\\
\begin{split}
\int_\Omega(\sigma_h)_{ij}(\tau_h)_{ij}\,dx
-\sum_{\kappa\in\mathcal T_h}
\int_\kappa\frac{\partial^2\psi_h}{\partial x_i\partial x_j}
(\tau_h)_{ij}\,dx\\
{}+\int_{\Gamma_h}M_n(\boldsymbol\tau_h)
S(\partial\psi_h/\partial n)\,ds
\end{split}
&=0
&&\forall\boldsymbol\tau_h\in\mathcal X_h.
\end{aligned}
\right.
\end{equation}
用 \eqref{eq:incompressible-hhj-form-a}
和 \eqref{eq:incompressible-hhj-form-b}, 可简写为
\begin{equation}\label{eq:incompressible-hhj-discrete-abstract}
\left\{
\begin{aligned}
b_h(\boldsymbol\sigma_h,\phi_h)
&=-\frac1\nu(\mathbf f,\operatorname{curl}\phi_h)
&&\forall\phi_h\in\Phi_h,\\
a_h(\boldsymbol\sigma_h,\boldsymbol\tau_h)
+b_h(\boldsymbol\tau_h,\psi_h)&=0
&&\forall\boldsymbol\tau_h\in\mathcal X_h.
\end{aligned}
\right.
\end{equation}

分析该问题之前, 必须仔细研究上述有限元空间,
特别是赋予适当范数并确定其自由度.
由 $b_h$ 的表达式, 取
\begin{equation}\label{eq:incompressible-hhj-tensor-mesh-norm}
\|\boldsymbol\tau\|_{0,h}
=\left(
\|\boldsymbol\tau\|_{0,\Omega}^2
+h\|M_n(\boldsymbol\tau)\|_{0,\Gamma_h}^2
\right)^{1/2}
\qquad\forall\boldsymbol\tau\in\mathcal X,
\end{equation}
以及
\begin{equation}\label{eq:incompressible-hhj-stream-mesh-norm}
\|\phi\|_{2,h}
=\left(
\sum_{\kappa\in\mathcal T_h}|\phi|_{2,\kappa}^2
+\frac1h\|S(\partial\phi/\partial n)\|_{0,\Gamma_h}^2
\right)^{1/2}
\end{equation}
对所有分片 $H^2$ 的 $\phi\in H^1(\Omega)$ 成立.
注意它与
\eqref{eq:incompressible-sfvp-mesh-norm-zero}
和 \eqref{eq:incompressible-sfvp-mesh-norm-two} 的类比.
还要注意
\[
|b_h(\boldsymbol\tau,\phi)|
\leq\|\boldsymbol\tau\|_{0,h}\|\phi\|_{2,h},
\qquad
|a_h(\boldsymbol\tau,\boldsymbol\mu)|
\leq\|\boldsymbol\tau\|_{0,\Omega}
\|\boldsymbol\mu\|_{0,\Omega}.
\]

\MathBookSourcePage{257}
现在确定 $\mathcal X_h$ 中张量的自由度.
先考虑 $\kappa$ 是参考三角形 $\widehat\kappa$ 的特殊情形.

\begin{Lemma}\label{lem:incompressible-hhj-reference-dofs}
在参考单位三角形 $\widehat\kappa$ 上,
对称张量值函数 $\boldsymbol\tau\in P_{l-1}^4$
由下列矩唯一确定:
\begin{equation}\label{eq:incompressible-hhj-reference-dofs}
\left\{
\begin{aligned}
&\int_{\widehat\kappa'}M_n(\boldsymbol\tau)q\,ds,
&&\forall q\in P_{l-1}(\widehat\kappa'),
\quad\text{对 }\partial\widehat\kappa\text{ 的每条边 }\widehat\kappa',\\
&\int_{\widehat\kappa}\tau_{ij}q\,d\widehat x,
&&\forall q\in P_{l-2}(\widehat\kappa),
\quad1\leq i\leq j\leq2.
\end{aligned}
\right.
\end{equation}
\end{Lemma}

\begin{Proof}
首先, \eqref{eq:incompressible-hhj-reference-dofs}
包含
\[
3l+\frac32l(l-1)=\frac32l(l+1)
\]
个自由度, 而
\[
\dim P_{l-1}^4=3\dim P_{l-1}
=\frac32l(l+1).
\]
因此只需证明相应齐次方程组
\begin{equation}\label{eq:incompressible-hhj-reference-homogeneous-dofs}
\left\{
\begin{aligned}
\int_{\widehat\kappa'}M_n(\boldsymbol\tau)q\,ds&=0
&&\forall q\in P_{l-1}(\widehat\kappa'),
\quad\forall\widehat\kappa'\subset\partial\widehat\kappa,\\
\int_{\widehat\kappa}\tau_{ij}q\,d\widehat x&=0
&&\forall q\in P_{l-2}(\widehat\kappa),
\quad1\leq i\leq j\leq2
\end{aligned}
\right.
\end{equation}
只有零解.

第一组方程等价于
$M_n(\boldsymbol\tau)=0$ 于 $\partial\widehat\kappa$ 上.
利用参考三角形的位置和形状, 即
\[
\tau_{22}=0\quad\text{on }\widehat x_2=0,
\qquad
\tau_{11}=0\quad\text{on }\widehat x_1=0,
\]
以及
\[
\tau_{11}+2\tau_{12}+\tau_{22}=0
\quad\text{on }\widehat x_1+\widehat x_2=1.
\]
在 \eqref{eq:incompressible-hhj-reference-homogeneous-dofs}
第二组方程中取 $q=\partial\tau_{ij}/\partial x_k$,
结合上述三个等式, 先得到
$\boldsymbol\tau=0$ 于 $x_1+x_2=1$, 再得到
$\boldsymbol\tau=0$ 于整个 $\partial\widehat\kappa$.
于是齐次矩条件立即蕴含
$\boldsymbol\tau=0$ 于整个 $\widehat\kappa$.
\end{Proof}

为把 Lemma~\ref{lem:incompressible-hhj-reference-dofs}
推广到任意三角形 $\kappa$, 必须引入适当变换,
把 $\widehat\kappa$ 上的对称张量映到 $\kappa$ 上的对称张量,
同时在某种意义下保持 $M_n(\cdot)$.
回顾 $\kappa$ 的法向量 $\mathbf n$
和 $\widehat\kappa$ 的法向量 $\widehat{\mathbf n}$ 满足
\[
\widehat{\mathbf n}
=\frac{B_\kappa^T(\mathbf n\circ F_\kappa)}
{\|B_\kappa^T(\mathbf n\circ F_\kappa)\|},
\qquad
\mathbf n\circ F_\kappa
=\frac{B_\kappa^{-T}\widehat{\mathbf n}}
{\|B_\kappa^{-T}\widehat{\mathbf n}\|},
\]
其中 $F_\kappa(\widehat x)=B_\kappa\widehat x+b_\kappa$,
$B_\kappa$ 是常系数非奇异矩阵.
另一方面,
\[
M_n(\boldsymbol\tau)(x)
=((\boldsymbol\tau(x)\mathbf n(x),\mathbf n(x))),
\]
这里 $((\cdot,\cdot))$ 是 $\mathbb R^2$ 的 Euclidean 内积.

\MathBookSourcePage{258}
由此在每个 $\kappa$ 上建立张量值函数之间的对应
\begin{equation}\label{eq:incompressible-hhj-tensor-piola-forward}
\widehat{\boldsymbol\tau}
=B_\kappa^{-1}(\boldsymbol\tau\circ F_\kappa)B_\kappa^{-T},
\end{equation}
等价地
\begin{equation}\label{eq:incompressible-hhj-tensor-piola-backward}
\boldsymbol\tau
=B_\kappa(\widehat{\boldsymbol\tau}\circ F_\kappa^{-1})B_\kappa^T.
\end{equation}
第一式显式写成
\[
\widehat\tau_{ij}
=\sum_{r,s=1}^2
(\tau_{rs}\circ F_\kappa)
\frac{\partial\widehat x_i}{\partial x_r}
\frac{\partial\widehat x_j}{\partial x_s}.
\]
显然该变换保持张量的对称性和正则性, 并且
\begin{equation}\label{eq:incompressible-hhj-normal-transform-forward}
M_{\widehat n}(\widehat{\boldsymbol\tau})
=\|B_\kappa^{-T}\widehat{\mathbf n}\|^2
(M_n(\boldsymbol\tau)\circ F_\kappa),
\end{equation}
或
\begin{equation}\label{eq:incompressible-hhj-normal-transform-backward}
M_n(\boldsymbol\tau)
=\|B_\kappa^T(\mathbf n\circ F_\kappa)\|^2
M_{\widehat n}(\widehat{\boldsymbol\tau}),
\end{equation}
因为
$\|B_\kappa^T(\mathbf n\circ F_\kappa)\|
=\|B_\kappa^{-T}\widehat{\mathbf n}\|^{-1}$.

因此容易验证 Lemma~\ref{lem:incompressible-hhj-reference-dofs}
可推广到 $\mathcal T_h$ 的任意三角形 $\kappa$.
从而可取下列量为 $\mathcal X_h$ 中张量的自由度:
\[
\int_{\kappa'}M_n(\boldsymbol\tau)q\,ds
\quad\forall q\in P_{l-1}(\kappa'),
\quad\forall\kappa'\subset\Gamma_h,
\]
以及
\[
\int_\kappa\tau_{ij}q\,dx
\quad\forall q\in P_{l-2}(\kappa),
\quad1\leq i\leq j\leq2,
\quad\forall\kappa\in\mathcal T_h.
\]

\begin{Remark}\label{rem:incompressible-hhj-linear-tensor-space}
$\mathcal X_h$ 最简单的例子对应 $l=1$:
\[
\mathcal X_h
=\{\boldsymbol\tau\in L^2(\Omega)^4;
\ \tau_{ij}=\tau_{ji},
\ \tau_{ij}|_\kappa\text{ 为常数},
\ \forall\kappa\in\mathcal T_h,
\ M_n(\boldsymbol\tau)\text{ 在 }\Gamma_h\text{ 每条边上连续}\}.
\]
显然, $\mathcal X_h$ 中的张量由每条边上的数据
$M_n(\boldsymbol\tau)$ 唯一确定.
\end{Remark}

\MathBookSourcePage{259}
同时可方便地定义一个从 $\mathcal X$ 到 $\mathcal X_h$
且具有良好性质的限制算子 $\Pi_h$.
先对 $\boldsymbol\tau\in L^2(\kappa)^4$,
$M_n(\boldsymbol\tau)\in L^1(\partial\kappa)$,
由下式定义 $P_{l-1}^4$ 中的对称张量
$\Pi_\kappa\boldsymbol\tau$:
\begin{equation}\label{eq:incompressible-hhj-local-restriction}
\left\{
\begin{aligned}
\int_{\kappa'}M_n(\Pi_\kappa\boldsymbol\tau
-\boldsymbol\tau)q\,ds&=0
&&\forall q\in P_{l-1}(\kappa'),
\quad\forall\kappa'\subset\partial\kappa,\\
\int_\kappa[(\Pi_\kappa\boldsymbol\tau)_{ij}-\tau_{ij}]q\,dx&=0
&&\forall q\in P_{l-2}(\kappa),
\quad1\leq i\leq j\leq2.
\end{aligned}
\right.
\end{equation}
再对 $\boldsymbol\tau\in\mathcal X$ 定义
\begin{equation}\label{eq:incompressible-hhj-global-restriction}
(\Pi_h\boldsymbol\tau)|_\kappa
=\Pi_\kappa(\boldsymbol\tau|_\kappa)
\qquad\forall\kappa\in\mathcal T_h.
\end{equation}
显然 $\Pi_h\boldsymbol\tau$ 对称,
而 $M_n(\boldsymbol\tau)$ 的连续性蕴含
$M_n(\Pi_h\boldsymbol\tau)$ 跨每条单元间边界连续,
所以 $\Pi_h\boldsymbol\tau\in\mathcal X_h$.

由 \eqref{eq:incompressible-hhj-tensor-piola-forward}
和 \eqref{eq:incompressible-hhj-normal-transform-forward},
\eqref{eq:incompressible-hhj-local-restriction}
等价于参考单元上的对应矩条件.
因此局部限制算子被仿射变换保持:
\begin{equation}\label{eq:incompressible-hhj-restriction-affine-commuting}
\widehat{\Pi_\kappa\boldsymbol\tau}
=\widehat\Pi(\widehat{\boldsymbol\tau})
\qquad\forall\kappa\in\mathcal T_h.
\end{equation}
其余性质由下一个 Lemma 给出.

\begin{Lemma}\label{lem:incompressible-hhj-restriction-properties}
假设剖分 $\mathcal T_h$ 正则.
由 \eqref{eq:incompressible-hhj-local-restriction}
和 \eqref{eq:incompressible-hhj-global-restriction}
定义的 $\Pi_h$ 是从 $\mathcal X$ 到 $\mathcal X_h$ 的线性映射, 且
\begin{equation}\label{eq:incompressible-hhj-restriction-orthogonality}
b_h(\boldsymbol\tau-\Pi_h\boldsymbol\tau,\phi_h)=0
\qquad\forall\phi_h\in\Phi_h,
\quad\forall\boldsymbol\tau\in\mathcal X,
\end{equation}
\begin{equation}\label{eq:incompressible-hhj-restriction-stability}
\|\Pi_h\boldsymbol\tau\|_{0,h}
\leq C_1\|\boldsymbol\tau\|_{0,h}
\qquad\forall\boldsymbol\tau\in\mathcal X.
\end{equation}
此外, 若某个实数 $k\in[1,l]$ 满足
$\boldsymbol\tau\in H^k(\Omega)^4\cap\mathcal X$, 则
\begin{equation}\label{eq:incompressible-hhj-restriction-error}
\|\boldsymbol\tau-\Pi_h\boldsymbol\tau\|_{0,h}
\leq C_2h^k|\boldsymbol\tau|_{k,\Omega}.
\end{equation}
\end{Lemma}

\begin{Proof}
由 \eqref{eq:incompressible-hhj-local-restriction},
$\Pi_h$ 显然是从 $\mathcal X$ 到 $\mathcal X_h$ 的线性映射.
又由同一式、$b_h$ 的表达式以及 $\Phi_h$ 的定义,
立即得到
\eqref{eq:incompressible-hhj-restriction-orthogonality}.

证明稳定性时, 首先由局部矩条件得到
\begin{equation}\label{eq:incompressible-hhj-normal-restriction-stability}
\|M_n(\Pi_h\boldsymbol\tau)\|_{0,\Gamma_h}
\leq\|M_n(\boldsymbol\tau)\|_{0,\Gamma_h}.
\end{equation}

\MathBookSourcePage{260}
再证明存在与 $h$ 无关的常数 $C>0$, 使
\begin{equation}\label{eq:incompressible-hhj-l2-restriction-stability}
\|\Pi_h\boldsymbol\tau\|_{0,\Omega}
\leq C\|\boldsymbol\tau\|_{0,h}.
\end{equation}
对每个 $\kappa\in\mathcal T_h$,
\eqref{eq:incompressible-hhj-restriction-affine-commuting}
和 \eqref{eq:incompressible-hhj-tensor-piola-forward} 给出
\[
\|\Pi_\kappa\boldsymbol\tau\|_{0,\kappa}^2
\leq|\det B_\kappa|\|B_\kappa\|^4
\|\widehat\Pi\widehat{\boldsymbol\tau}\|_{0,\widehat\kappa}^2.
\]
由参考单元上的矩条件,
\[
\|\widehat\Pi\widehat{\boldsymbol\tau}\|_{0,\widehat\kappa}
\leq C_1\left(
\|\widehat{\boldsymbol\tau}\|_{0,\widehat\kappa}
+\|M_{\widehat n}(\widehat{\boldsymbol\tau})\|_{0,\partial\widehat\kappa}
\right).
\]
利用法向量变换、逆变换和尺度关系, 汇总后得到
\[
\|\Pi_\kappa\boldsymbol\tau\|_{0,\kappa}^2
\leq C_2\sigma^6\left(
\|\boldsymbol\tau\|_{0,\kappa}^2
+h_\kappa\|M_n(\boldsymbol\tau)\|_{0,\partial\kappa}^2
\right).
\]
剖分正则性于是证明
\eqref{eq:incompressible-hhj-l2-restriction-stability},
进而结合
\eqref{eq:incompressible-hhj-normal-restriction-stability}
得到 \eqref{eq:incompressible-hhj-restriction-stability}.

最后证明误差估计.
同样有
\[
\|\boldsymbol\tau-\Pi_h\boldsymbol\tau\|_{0,\kappa}^2
\leq|\det B_\kappa|\|B_\kappa\|^4
\|\widehat{\boldsymbol\tau}
-\widehat\Pi\widehat{\boldsymbol\tau}\|_{0,\widehat\kappa}^2.
\]
由于 $\widehat\Pi$ 保持系数属于 $P_{l-1}$ 的对称张量,
Bramble--Hilbert 估计给出
\[
\|\widehat{\boldsymbol\tau}
-\widehat\Pi\widehat{\boldsymbol\tau}\|_{0,\widehat\kappa}
\leq C_3|\widehat{\boldsymbol\tau}|_{k,\widehat\kappa}.
\]
由 \eqref{eq:incompressible-hhj-tensor-piola-backward}
和尺度关系,
\[
|\widehat{\boldsymbol\tau}|_{k,\widehat\kappa}
\leq C_5\|B_\kappa^{-1}\|^2
|\det B_\kappa^{-1}|^{1/2}
h_\kappa^k|\boldsymbol\tau|_{k,\kappa}.
\]
剖分正则性于是给出
\[
\|\boldsymbol\tau-\Pi_h\boldsymbol\tau\|_{0,\Omega}
\leq C_6h^k|\boldsymbol\tau|_{k,\Omega}.
\]
类似地, 由
\eqref{eq:incompressible-hhj-normal-transform-backward}
和 trace Theorem 得
\[
\|M_n(\boldsymbol\tau-\Pi_h\boldsymbol\tau)\|_{0,\partial\kappa}^2
\leq C_9h^{2k-1}|\boldsymbol\tau|_{k,\kappa}^2.
\]
从而得到
\eqref{eq:incompressible-hhj-restriction-error}.
\end{Proof}
